%\setlength{\columnsep}{12pt}      % espaço entre colunas
%\setlength{\columnseprule}{0.2pt} % linha vertical entre colunas


%-------------------------------------------------------------------------
% Modelo Trabalho PPGE-UFPB/Estatística
% Normas da ABNT
% Provedor: Laboratório de Economia e Modelagem Aplicada -- LEMA
%-------------------------------------------------------------------------

\documentclass[12pt, openright,oneside, a4paper, english, brazil]{report}

% ----------------------------
% Pacotes básicos 
% ----------------------------
% Fonte e codificação (acentuação)
\usepackage{lmodern}

\renewcommand{\sfdefault}{ppl}
\usepackage[T1]{fontenc}	
\usepackage[utf8]{inputenc}	
\usepackage[brazil]{babel}
\usepackage{comment} % para esconder a tabela Org do compilador
% Referências
%\usepackage[brazilian,hyperpageref]{}	
%\usepackage[alf]{abntex2cite}			
\usepackage[colon]{natbib}
\bibliographystyle{apalike}

%Texto aleatório
\usepackage{lipsum} 

% Tabelas e Figuras
\usepackage{ctable}             
\usepackage{multirow}
\usepackage{array}
\usepackage{float}              
\usepackage{longtable}          
\usepackage{xcolor,colortbl}    
\usepackage{booktabs}           
\usepackage{graphicx}			
\usepackage{subfig}             
\usepackage{caption}            
\usepackage{epstopdf}           

% Equações e símbolos
\usepackage{amsmath}            
\usepackage{nomencl} 
\usepackage{hyperref, cleveref}  

% Texto
\usepackage{indentfirst}		
\usepackage{color}				
\usepackage{microtype} 			
\usepackage{lastpage}			
% ----------------------------------------------------------
% Configurações do texto
% ----------------------------------------------------------
\usepackage{multicol}
\usepackage{array}
\newcolumntype{L}[1]{>{\raggedright\let\newline\\\arraybackslash\hspace{0pt}}m{#1}}
\newcolumntype{C}[1]{>{\centering\let\newline\\\arraybackslash\hspace{0pt}}m{#1}}
\newcolumntype{R}[1]{>{\raggedleft\let\newline\\\arraybackslash\hspace{0pt}}m{#1}}

\makeatletter
%\hypersetup{pdftitle={\@title}, pdfauthor={\@author}}
\makeatother

\def\preambulo{Trabalho elaborado em cumprimento às exigências da disciplina \disciplina, ministrada pela professora Dra. \professor.}
%--------------------------------------------------------------

\setlength{\parindent}{1.25cm}
\setlength{\parskip}{0.1cm} 

\usepackage{titlesec}
\usepackage{lipsum}

\makeindex





\title{Atividade 01 - Ética}

\def\autor{Oscar Sérgio Monteiro\\
\texttt{oscar.monteiro@ufpr.br}}

\def\disciplina{SC325 - Ética Geral e Profissional, turma a}
\def\professor{Edicreia Santos}


\begin{document}
\makeatletter
\begin{titlepage}
    \begin{figure}[H]
    \vspace{-3cm}
    \centering
    \includegraphics[height=3cm,keepaspectratio]{config/coverpage}\hfill
    \par\vspace{1cm}
    \end{figure}
    \begin{figure}[H]
    \vspace{-2cm}
    \centering
    \includegraphics[height=3cm,keepaspectratio]{config/contabeis.png}
    \par\vspace{1cm}
    \end{figure}
    
    \vspace{2cm}
    
	\begin{center}
	    {\huge \bfseries \@title}
	    \begin{flushright}
		\begin{minipage}{8cm}
			\preambulo
		\end{minipage}
	\end{flushright}
	\vspace{3cm}
	    \normalsize
	    \autor
	\end{center}
	
\vfill
\centering
{\large Curitiba-PR, \the\year}

\end{titlepage}
\makeatother


\pdfbookmark[0]{\contentsname}{toc}
\tableofcontents
\newpage


\titleformat{\chapter}[display]{\normalfont\bfseries}{}{0pt}{\Huge  \thechapter. }
\newcommand{\doubletoprule}{%
  \specialrule{.1em}{2pt}{2pt}%
  \specialrule{.1em}{0pt}{0pt}%
}

\chapter{Estudo de caso}
\subsection{Descrição do casos}
Uma empresa brasileira de tecnologia desenvolveu um aplicativo que promete aumentar a produtividade das pessoas no trabalho realizado na modalidade de teletrabalho (fora dos espaços físicos da empresa), mas para isso, coleta dados pessoais detalhados dos usuários, como histórico de navegação e hábitos de consumo. A empresa desenvolvedora do aplicativo afirma que essas informações serão usadas apenas para melhorar a experiência do usuário e nunca serão compartilhadas. Contudo, isso não foi claramente informado no momento da instalação do aplicativo, e muitas pessoas não percebem que seus dados estão sendo coletados.\\

\subsection{item 1}
1 Se você fosse o CEO dessa empresa de tecnologia, o que faria? Você priorizaria o lucro da empresa e o potencial de crescimento ou garantiria a transparência sobre a coleta de dados e o consentimento explícito dos usuários, mesmo que isso resultasse em perda de clientes? Quais possíveis benefícios e consequências você consideraria para sua decisão?\\
\subsection{Posicionamento}
Não tenho dúvida que a melhor alternativa para a empresa é estar em conformidade com a LGPD (Lei Geral de Proteção de Dados) e solicitar o consentimento dos usuários (cliente) para o uso das informações que estão sendo coletadas. Essa é uma prática bem usual. De modo geral os usuários são displicentes em relação a privacidade de seus dados, preferem aplicativos baratos ou de uso gratuito e oferecem em troca seu potencial de compra.\\
\subsection{item 2}
2. Você trabalha na área contábil dessa empresa e, após descobrir como os dados pessoais dos usuários estão sendo coletados sem o devido consentimento, você decide vazar essas informações para a imprensa. Sua ação traria grande repercussão sobre a ética da empresa. Você tomaria essa atitude? Quais seriam as consequências para você e para a empresa?\\
\subsection{Posicionamento}
De forma alguma denunciaria a empresa a imprensa nem a Agência Nacional de Proteção de Dados. Eu reportaria o problema a meu superior direto. A segunda providência seria revisar os procedimentos do setor de contabilidade para garantir que as informações usadas na contabilidade estão seguras, se algum risco significativo de vazamento houver, encaminharia um plano de correção ao meu supervisor.  Procuraria uma formalização da posição da empresa. Caso tome conhecimento de que o fato chegou a direção e responsáveis e que a empresa decidiu assumir o risco da irregularidade, estaria em com conflito ético de pedir ou não demissão. Apenas se vivenciar esta questão de verdade poderia saber qual seria meu posicionamento, os dados descritos são muitos superficiais para uma decisão de demissão.
\\
Os riscos para a empresa são os de imagem e de multa. Para mim não existiriam riscos, não sou o responsável pela questão e agi de boa fé quando soube do problema.\\
\end{document}
